%! AP Statistics — Unit 7, Lesson 7.2 — Guided Notes (Answer Key)
\documentclass[10pt]{article}
\usepackage{apstats-article}
\usepackage{apstats-key}
\newcommand{\TallMath}[1]{$\displaystyle #1\rule[-1.4em]{0pt}{3.2em}$}

\begin{document}

\pageheader{Unit 7, Lesson 7.2}{Guided Notes --- Answer Key}
\namedateperiod

\begin{objectivebox}
By the end of this lesson, I will be able to:
\begin{itemize}
    \item[$\square$] \ansline{Describe the $t$-distribution and explain how it differs from the standard normal.}
    \item[$\square$] \ansline{Identify and verify the conditions for a one-sample $t$-interval.}
    \item[$\square$] \ansline{Calculate a one-sample $t$-interval: $\bar x \pm t^* \cdot s/\sqrt{n}$.}
    \item[$\square$] \ansline{Apply the one-sample $t$-interval to a matched-pairs design using differences.}
\end{itemize}
\end{objectivebox}

\begin{vocabbox}
\termblanklong{$t$-distribution}
\ansline{A symmetric, bell-shaped family of distributions with heavier tails than $N(0,1)$; used when $\sigma$ is unknown and $s$ is substituted; indexed by degrees of freedom.}

\termblanklong{Degrees of freedom ($df$)}
\ansline{For a one-sample $t$-procedure: $df = n - 1$.}

\termblanklong{Standard error of $\bar x$ ($SE_{\bar x}$)}
\ansline{$SE_{\bar x} = s/\sqrt{n}$; estimates the standard deviation of the sampling distribution of $\bar x$.}

\termblanklong{Critical value $t^*$}
\ansline{The positive value such that the area between $-t^*$ and $t^*$ under the $t$-curve equals $C$; found from the $t$-table with $df = n - 1$.}

\termblanklong{One-sample $t$-interval}
\ansline{$\bar x \pm t^* \cdot s/\sqrt{n}$; the appropriate CI for a population mean when $\sigma$ is unknown.}

\termblanklong{Matched-pairs design}
\ansline{A design where each observation in one group is paired with an observation in the other; analyzed by computing differences and applying a one-sample $t$-procedure.}
\end{vocabbox}

\begin{hookbox}
Why can't we use a $z$-interval for the sodium example?
\ansline{$\sigma$ (the true population standard deviation of sodium intake) is unknown. A $z$-interval requires $\sigma$; substituting $s$ changes the distribution of the standardized statistic from $z$ to $t$.}
\end{hookbox}

\begin{notesbox}{The $t$-Distribution (VAR-7.A)}
When we substitute $s$ for $\sigma$, the statistic $t = \dfrac{\bar x - \mu}{s/\sqrt{n}}$
follows a \ans{$t$}-distribution.

\textbf{Properties:}
\begin{itemize}
    \item Bell-shaped, \ans{symmetric}, centered at \ans{0}.
    \item Heavier tails than the standard normal.
    \item Degrees of freedom: $df = n - $ \ans{1}.
    \item As $df$ \ans{increases}, the $t$-distribution approaches $N(0, 1)$.
\end{itemize}

For 95\% CI with $n = 20$: $df = $ \ans{19}, $t^* = $ \ans{2.093}.
\end{notesbox}

\begin{notesbox}{Conditions for a One-Sample $t$-Interval (UNC-4.P)}
\begin{enumerate}
    \item \textbf{Random:} Data come from a \ans{random sample} or randomized experiment.
    \item \textbf{10\%:} $n \le $ \ans{10}\%$ \cdot N$ (if sampling w/o replacement).
    \item \textbf{Large Sample:}
          \begin{itemize}
              \item $n \ge $ \ans{30}: condition met.
              \item $n < $ \ans{30}: data show no strong \ans{skewness} or \ans{outliers}.
          \end{itemize}
\end{enumerate}
\end{notesbox}

\begin{notesbox}{The One-Sample $t$-Interval Formula (UNC-4.Q, UNC-4.R)}
\renewcommand{\arraystretch}{2}
\begin{tabularx}{\linewidth}{lX}
    Point estimate & $\bar x$ \\
    Standard error & $SE_{\bar x} = s/\sqrt{n}$ \\
    Margin of error & $ME = t^* \cdot s/\sqrt{n}$ \\
    \textbf{Interval} & $\bar x \pm t^* \cdot s/\sqrt{n}$ \\
\end{tabularx}

\textbf{Sodium example:} $df = $ \ans{19}, $t^* = $ \ans{2.093},
$SE = 480/\sqrt{20} = $ \ans{107.3} mg

Interval: $2840 \pm 2.093 \cdot 107.3 = 2840 \pm 224.6
= ($ \ans{2615} $,$ \ans{3065} $)$ mg
\end{notesbox}

\begin{notesbox}{The Four-Step Procedure (PCCF)}
\begin{enumerate}
    \item \textbf{Parameter:} \ansline{$\mu = $ the true mean [quantity] for [population].}
    \item \textbf{Conditions:} Random, 10\%, Large Sample --- all verified.
    \item \textbf{Calculate:} $\bar x \pm t^* \cdot s/\sqrt{n}$
    \item \textbf{Conclude:} ``We are \ans{$C$}\% confident that the interval captures
          the true mean \ans{[quantity]} for \ans{[population]}.''
\end{enumerate}
\end{notesbox}

\begin{notesbox}{Matched-Pairs Design (UNC-4.O.3)}
\begin{enumerate}
    \item Compute $d_i = x_{i,1} - x_{i,2}$.
          Order of subtraction: \ansline{Define clearly (e.g., after $-$ before), and use consistently.}
    \item $\bar d = $ \ans{mean of the differences} \quad $s_d = $ \ans{SD of the differences}.
    \item Interval: $\bar d \pm t^* \cdot s_d/\sqrt{n}$, $df = n - 1$,
          where $n = $ number of \ans{pairs}.
\end{enumerate}
\end{notesbox}

\end{document}
